\section{Conclusion}

% - 呼应关键负结果（不要淡化）：把VCWorld的检索增强LLM因果推理框架移植到酵母后，
%   "扰动->表达"（T1）复现了VCWorld在人类数据上的规模趋势，但"扰动->生长表型"（T2A）
%   没有——四个模型在T2A-DIR上都接近随机，且不随模型规模提升。这是本文的核心发现，
%   结论段应该正面陈述它，而不是把负结果包装成"未来工作"淡化掉。
% - 对"虚拟细胞"建模范式的启示（呼应RQ1-4的发现，具体措辞待RQ4机制一致率/RQ3消融数字
%   定稿后再收紧）：纯文本知识检索+LLM推理这条白盒路线，长于"扰动->调控/表达"这类知识
%   图谱能稠密编码因果链（TF->靶基因符号）的问题；短于"扰动->定量剂量依赖的适合度"这类
%   需要连续数值信号（浓度、通路负荷、代谢流量）而非离散关系型知识的问题。对未来虚拟细胞
%   框架的实践建议：不同读出层可能需要混合架构——调控类读出继续用白盒检索+推理，适合度/
%   定量类读出补充数值特征或轻量监督模块（呼应 RQ3 训练型ML对照里GAT/KGE-TGI等在T1上
%   有时甚至超过LLM的观察，见docs/task_results.md）。
% - Limitations：（1）酵母单一物种，人类->酵母的知识密度差异本身是否是T1/T2差距的混杂
%   因素，需要在人类T2类任务上做对照才能完全排除；（2）T2A目前只用Hillenmeyer 2008单一
%   数据集，T2B(SGA遗传互作)、het/HIP等平行任务尚未接入（docs/task_definitions.md"二、
%   可能的任务"），负结果的普适性有待更多生长表型数据源验证；（3）本地开源模型规模上限
%   14B，未接入更强的闭源模型（Gemini等，PLAN.md决策记录4"无API"），T1/T2差距在更强模型
%   上是否收窄尚未知；（4）RQ4可解释性核查、部分消融/训练型对照写作时状态需要先核实
%   （见Experiments节内联注释），结论段的具体措辞要等这些数字定稿后收紧，避免过度断言。
